\section{\uppercase{Results}}

\subsection{Per-Modality Encoder Evaluation}

\noindent
Before the two streams are fused, each encoder is scored against the standard
benchmark of its field, so that any fusion result rests on encoders already
tested in isolation. Table~\ref{tab:perleg} collects both comparisons.

On UJIIndoorLoc~\cite{torressospedra2014ujiindoorloc}, WiFi-Net reaches an MAE
of $8.58$~m on the validation scans, against $15.17$~m for wlan\_localization~%
\cite{wlanloc}, a reduction of $43.4\%$ under Eq.~\eqref{eq:mae}.
% src: experiments-results-brief.md 2.1; task grounded number 8.58 m / -43.4%
The benchmark fixes $520$ access points and carries no trajectory structure, so
the comparison isolates the fingerprint-to-position mapping with no temporal
context to lean on. Under that constraint WiFi-Net more than halves the
reference error.

The inertial leg is the harder case. On RoNIN canonical~\cite{yan2019ronin},
IMUCNN attains a raw ATE of $9.72$~m on the $32$ held-out sequences, against
$5.14$~m for ResNet1D~\cite{yan2019ronin} under Eq.~\eqref{eq:ate}, so the
lighter encoder trails the reference by roughly $89\%$.
% src: task grounded number raw ATE 9.72 m vs 5.14 m, ~+89% gap
The two encoders are not the same size: IMUCNN carries about $0.05$~M
parameters against ResNet1D's $4.6$~M, a factor of roughly $95$. The gap
reflects an in-domain-competitive encoder paying a cross-subject generalization
cost at a far smaller parameter budget, rather than a failure of the encoder on
held-out motion. Figure~\ref{fig:ronin_traj} shows the qualitative consequence
on one held-out sequence.

\begin{table*}[tb]
\caption{Per-modality encoder evaluation, each encoder scored on the standard
  benchmark of its field (MAE on UJIIndoorLoc, raw ATE on RoNIN canonical).
  Bold marks the lower error in each row.}\label{tab:perleg}
\centering
\begin{tabular}{|l|l|c|c|c|}\hline
Benchmark (split) & Metric & wlan\_localization / ResNet1D & Ours & $\Delta$ \\ \hline
UJIIndoorLoc (val) & MAE & 15.17 & \textbf{8.58} & $-43.4\%$ \\ \hline
RoNIN canonical & raw ATE & \textbf{5.14} & 9.72 & $+89\%$ \\ \hline
\end{tabular}
\end{table*}

\begin{figure}[!hb]\centering
  \includegraphics[width=0.95\columnwidth]{figures/fig_ronin_traj.pdf}
  \caption{RoNIN canonical held-out sequence: ground truth, ResNet1D, and
    IMUCNN (ours). The lighter inertial encoder follows the route but
    accumulates more drift than ResNet1D over the leg.}\label{fig:ronin_traj}
  % [[VERIFY: export figure without in-plot title]]
\end{figure}

\subsection{End-to-End Fusion}

\noindent
With both encoders evaluated, the set-transformer fuses them into a single
position estimate. The cross-session real-world collection carries the main
claim; the controlled simulation is reported first as a sanity check.
Table~\ref{tab:fusion} collects the fusion results against the named baselines.

MSILN site1/B1 is the cross-session test: training, validation and test come
from sessions recorded on separate days. On the held-out test session the
set-transformer reaches $10.89$~m MAE under Eq.~\eqref{eq:mae}, against
$28.31$~m for wlan\_localization, an improvement of $61.5\%$.
% src: task grounded numbers MSILN test 10.89 m vs wlanloc 28.31 m, -61.5%
On the validation session the proposed method reports $15.22$~m against
$21.26$~m for wlan\_localization, a smaller margin of $28.4\%$.
% src: task grounded numbers MSILN val 15.22 m vs wlanloc 21.26 m, -28.4%
Where the WiFi baseline nearly doubles its error from validation to test, the
fusion error falls, because the dense inertial stream holds position against the
sparse WiFi anchor across the session boundary. The proposed method therefore
improves over the open-source WiFi baseline under genuine cross-session
fingerprint drift, which is the setting it is built for. Figure~\ref{fig:msiln_cdf}
reports the full per-sample error distribution on the test session.

This margin should be read with one caveat. A WiFi-kNN comparator is comparable
to the proposed method on this test split, because a single WiFi-dense path
contributes roughly $28\%$ of the test mass and is easy for nearest-neighbour
fingerprinting (Section~\ref{sec:limitations}).
% src: scope.md 7.3 path-130 ~28% of test mass
The cross-session improvement is therefore reported against the open-source
wlan\_localization baseline, not as a claim of dominance over every WiFi method.

In the Webots simulation, the two-modality model at $K=4$ reaches about
$0.375$~m MAE on test.
% src: task grounded number Webots 2-mod K=4 ~0.375 m
% [[VERIFY: Webots canonical run -- val/test exact figure unresolved, 0.375 vs 0.517]]
No public two-modality method targets this simulation, and the WiFi field is
synthesized by Gaussian-process regression rather than measured, so the
sub-metre figure is an optimistic controlled-simulation check rather than a
real-deployment result. It is reported only to confirm that the pipeline
recovers position when the WiFi field is clean.

\begin{table*}[tb]
\caption{End-to-end fusion: position MAE (m) under Eq.~\eqref{eq:mae}. The
  simulation has no public two-modality reference and uses a synthesized WiFi
  field, so its sub-metre figure is a controlled-simulation check, not a
  real-deployment result. The MSILN comparison is cross-session against the
  open-source WiFi baseline. Bold marks the lowest error in each row.}%
  \label{tab:fusion}
\centering
\begin{tabular}{|l|c|c|c|}\hline
Dataset (split) & wlan\_localization & Ours & $\Delta$ \\ \hline
Webots sim (test) & -- & \textbf{$\sim$0.375} & -- \\ \hline
MSILN site1/B1 (val) & 21.26 & \textbf{15.22} & $-28.4\%$ \\ \hline
MSILN site1/B1 (test) & 28.31 & \textbf{10.89} & $-61.5\%$ \\ \hline
\end{tabular}
\end{table*}

\begin{figure}[!hb]\centering
  \includegraphics[width=0.95\columnwidth]{figures/fig_msiln_cdf.pdf}
  \caption{Per-sample position error on the MSILN site1/B1 cross-session test
    session: cumulative distribution for the proposed method against
    wlan\_localization. The proposed method places most samples within a
    shorter error range, while the WiFi baseline carries a long tail under
    cross-session drift.}\label{fig:msiln_cdf}
  % [[VERIFY: export figure without in-plot title]]
\end{figure}

The IMUWiFine LSTM~\cite{nurpeiissov2022imuwifine}, a learned WiFi--IMU fusion,
is the natural learned-fusion comparator for the cross-session setting. No
grounded number for the proposed method against it on IMUWiFine floor~4 is
available at the time of writing, so that comparison is deferred.
% [[VERIFY: our method vs IMUWiFine-LSTM (nurpeiissov2022imuwifine) on IMUWiFine fl.4 numbers needed]]

\subsection{Robustness}

\noindent
The robustness experiments run on the Webots simulation, where every signal is
controlled and individual sensors can be suppressed or aged on demand. The
sub-metre scale of these numbers carries the synthesized-WiFi caveat noted
above; the evidence here is the shape of the degradation, not its absolute
level.

\textbf{Modality ablation.} Suppressing a modality at test time exposes what
each stream contributes. With WiFi only the test MAE is about $0.495$~m; with
IMU only it is about $3.77$~m; with both it is about $0.375$~m.
% src: task grounded numbers wifi-only ~0.495, imu-only ~3.77, both ~0.375
The WiFi token supplies the absolute anchor, so dropping IMU costs little, while
IMU alone drifts without an absolute reference. The model degrades gracefully
rather than failing when one stream is absent.

\textbf{Staleness.} Aging the WiFi token while the IMU stream stays fresh maps
the cost of asynchronous, delayed measurements. The test error rises smoothly
with WiFi lag at a slope of about $0.029$~m/s, rather than collapsing once the
anchor goes stale.
% src: task grounded number staleness slope ~0.029 m/s
Small lags cost little, so what would be a cliff under a fixed-rate design
becomes a slope across moderate delays.

\textbf{Latency.} On the Quadro P4000 the model runs at about $5.1$~ms per
sample at batch size $1$ and about $0.17$~ms per sample at batch size $32$,
fast enough for real-time use.
% src: task grounded numbers latency ~5.1 ms (b=1), ~0.17 ms (b=32)

\subsection{Limitations}\label{sec:limitations}

\noindent
The inertial encoder trails the reference on its own benchmark. On RoNIN
canonical IMUCNN reports a raw ATE of $9.72$~m against ResNet1D's $5.14$~m
($+89\%$); removing a rigid-body offset by Umeyama alignment narrows the gap to
roughly $7.6$--$7.9$~m but does not close it.
% src: task -- Umeyama-aligned ATE ~7.6-7.9 m belongs only to limitations
Even with alignment the gap remains, and it reflects a $95\times$-smaller
encoder paying a cross-subject generalization cost. Closing it at this parameter
budget is open work.

The fused trajectory is over-smoothed relative to ground truth, with predicted
motion magnitude falling below the true magnitude at turns. The temporal
attention learns to smooth, and at this loss setting it smooths too much.
Despite this, absolute-position accuracy stays competitive on the cross-session
test, and an auxiliary velocity loss is identified as a candidate fix.

The MSILN test split is not uniform. One path contributes roughly $28\%$ of the
test mass and is WiFi-dense, so its WiFi-favourable composition shapes the
aggregate test figure and explains why a WiFi-kNN comparator is comparable on
this split. The cross-session improvement reported above holds against the
open-source WiFi baseline, but the test figure should be read alongside the
validation figure, not in isolation.
